% Cross-document references synchronized for the 31 October 2026 revision.
\documentclass[11pt]{article}
\usepackage[T1]{fontenc}
\usepackage[utf8]{inputenc}
\usepackage{lmodern,microtype}
\usepackage[letterpaper,margin=1.08in]{geometry}
\usepackage{amsmath,amssymb,amsthm,mathtools}
\usepackage{enumitem,booktabs,array}
\usepackage[round,authoryear]{natbib}
\usepackage{xcolor}
\usepackage[colorlinks=true,linkcolor=blue!40!black,citecolor=blue!40!black,urlcolor=blue!40!black]{hyperref}
\usepackage{fancyhdr}
\pagestyle{fancy}\fancyhf{}\fancyhead[L]{\small Definite Extensions: Technical Supplement}\fancyhead[R]{\small\thepage}
\setlength{\headheight}{14pt}
\setlength{\parskip}{3pt}
\setlength{\emergencystretch}{2em}
\setlist{nosep}
\newtheorem{theorem}{Theorem}[section]
\newtheorem{proposition}[theorem]{Proposition}
\newtheorem{lemma}[theorem]{Lemma}
\newtheorem{corollary}[theorem]{Corollary}
\theoremstyle{definition}
\newtheorem{definition}[theorem]{Definition}
\newtheorem{remark}[theorem]{Remark}
\newcommand{\D}{\mathsf D}
\newcommand{\poss}{\mathsf M}
\newcommand{\Rig}{\operatorname{Rig}}
\newcommand{\Pers}{\operatorname{Pers}}
\newcommand{\Inext}{\operatorname{Inext}}
\newcommand{\ED}{\operatorname{ED}}
\newcommand{\Def}{\operatorname{Def}}
\newcommand{\ER}{\operatorname{ER}}
\newcommand{\ID}{\operatorname{ID}}
\newcommand{\Pl}{\operatorname{Pl}}
\newcommand{\Co}{\operatorname{Co}}
\newcommand{\Hull}{\operatorname{Hull}}
\newcommand{\lev}{\operatorname{lev}}
\newcommand{\Arr}{\operatorname{Arr}}
\newcommand{\tg}{\operatorname{tg}}
\newcommand{\id}{\mathrm{id}}
\newcommand{\C}{\mathsf C}
\newcommand{\HD}{\mathsf H_{\D}}
\newcommand{\U}{\mathbf U}
\newcommand{\Z}{\mathbf 0}
\newcommand{\sng}[1]{\delta_{#1}}
\newcommand{\restrall}[2]{\forall^{#1}x\,#2}
\newcommand{\rel}{\mathbin{\approx}}
\newcommand{\deq}{\mathbin{\equiv_{\D}}}
\newcommand{\sub}{\mathbin{\sqsubseteq}}
\newcommand{\ND}{\mathsf{ND}_e}

\numberwithin{equation}{section}
\title{Definite Extensions\\[3pt]\large Technical Supplement}
\author{AI-generated research draft}
\date{31 October 2026}

\begin{document}
\maketitle
This supplement contains the logical calculus, proofs, constructions, and countermodels for \emph{Definite Extensions: Rigidity, Circumscription, and Critical Plural Logic}, revised 31 October 2026. Main-text Sections 1--9 and unlettered equation numbers refer to that paper; sections and results A--F belong to this supplement. The central results are stated where they are used in the main text. The table identifies their proofs.

\begin{center}\small
\begin{tabular}{@{}p{0.14\textwidth}p{0.54\textwidth}p{0.23\textwidth}@{}}
\toprule
Main text & Claim & Proof or construction\\
\midrule
\S3 & Quantifier characterization; comparison under the basis condition & A.2; B.1--B.2\\
\S4 & Representation, identity, and the separation of definiteness notions & A.2; B.3; F.3\\
\S5 & Closure results and limits of the base theory & C; F.1\\
\S6 & Rigid-property classification, critical principles, independence of Infinity & D\\
\S7 & Set-theoretic consequences and relative consistency & E\\
\S8 & Constant-domain and S5 obstructions & F.2\\
\bottomrule
\end{tabular}\end{center}

\subsection*{Notation and assumptions}
Individuals have type $e$ and propositions type $t$. Relational types are generated by $\sigma\to\tau$ with $\tau\ne e$. The typed calculus $\C$, and modal fragment $\HD$ are specified in Section~\ref{app:logic}. Broad necessity is $\Box p:=(p=\top)$; $\D:t\to t$ is a genuine definiteness operator, with dual $\poss p:=\neg\D\neg p$. Write
\begin{align*}
F\rel G&:=\forall x(Fx\leftrightarrow Gx), &
F\sub G&:=\forall x(Fx\rightarrow Gx),\\
F\deq G&:=\D(F\rel G), &
\forall^R x\varphi&:=\forall x(Rx\rightarrow\varphi).
\end{align*}
For $N=\D$ or $N=\Box$, the definitions are
\begin{align*}
\Pers_N(R)&:=N\forall x(Rx\rightarrow NRx),\\
\Inext_N(R)&:=N\forall Q(\forall^R x\,NQx\rightarrow N\forall^R x\,Qx),\\
\Rig_N(R)&:=\Pers_N(R)\wedge\Inext_N(R),\\
\Rig(R)&:=\Rig_{\Box}(R),\qquad \Def(R):=\Rig_{\D}(R),\\
\ER(F)&:=\exists R(\Rig(R)\wedge F\rel R),\\
\ED(F)&:=\exists R(\Def(R)\wedge F\rel R),\\
\Pl(R)&:=\Rig(R)\wedge\exists xRx.
\end{align*}
Also $\Z:=\lambda x.\bot$, $\U:=\lambda x.x=x$, and $\sng a:=\lambda x.x=a$. Pointwise definiteness is $\ID_\D(F):=\forall x(\D Fx\vee\D\neg Fx)$; necessary individual distinctness is $\ND:=\D\forall xy(x\ne y\rightarrow\D(x\ne y))$.

The comparison theorem assumes, throughout broad possibility, normal S4 and converse Barcan for $\D$ at every type, together with
\[
\mathsf{MB}:\quad\Box\exists K^{t\to t}\forall p
\bigl(\D p\leftrightarrow\forall q(Kq\rightarrow\Box(q\rightarrow p))\bigr).
\]
The critical theory $\mathsf{CP}$ adds broad Generalized Union and Infinity to $\C$, with their broad necessitations. Writing $B_x:=\lambda y.Bxy$, these are
\[
\ER(A)\wedge\forall^A x\ER(B_x)\rightarrow
\ER(\lambda y\exists x(Ax\wedge Bxy)),
\]
\[
\begin{aligned}
\forall x\exists!y\psi(x,y)\ \longrightarrow\
\forall R\bigl(\Rig(R)\rightarrow\exists S[\Rig(S)\wedge R\sub S\wedge{}\\
\forall xy(Sx\wedge\psi(x,y)\rightarrow Sy)]\bigr).
\end{aligned}
\]
The functional relation is evaluated currently. The combined set theory $\mathsf{CS}$ adds the stated $\D$ assumptions and $\mathsf{MB}$, a set predicate $\mathsf S$, membership $\in$, and the following axioms and their broad necessitations. With $E_s:=\lambda x.x\in s$:
\begin{align*}
&\forall xs(x\in s\rightarrow\mathsf Ss),\\
&\forall st(\mathsf Ss\wedge\mathsf St\wedge E_s\rel E_t\rightarrow s=t),\\
&\forall s(\mathsf Ss\rightarrow\ER(E_s)),\\
&\forall F(\ED(F)\rightarrow\exists s(\mathsf Ss\wedge E_s\rel F)).
\end{align*}
These are typing, Extensionality, set-boundedness, and Definite Set Formation. Their consequences and joint interpretation are proved in Section~\ref{app:sets}.
\clearpage
\appendix

\section{The logical calculus and elementary representation results}\label{app:logic}

\subsection{The full higher-order background}

Use the relational types specified in Section~3. The logical constants are classical propositional connectives, identity at each type, and quantifiers at each type. Take $\exists v\varphi$ as $\neg\forall v\neg\varphi$ if desired. A precise background calculus $\mathsf H$ consists of classical tautologies and modus ponens, universal generalization, and the schemas
$$
\begin{gathered}
\forall v\varphi\rightarrow\varphi[A/v],\qquad
\forall v(\varphi\rightarrow\psi)\rightarrow(\varphi\rightarrow\forall v\psi)
   \quad(v\text{ not free in }\varphi),\\
A=A,\qquad A=B\rightarrow(FA\rightarrow FB),\\
(\lambda v.M)A=M[A/v],\qquad \lambda v.Fv=F
   \quad(v\text{ not free in }F).
\end{gathered}
$$
All expressions are well typed, substitution is capture-avoiding, and generalization has the usual restriction on undischarged assumptions. The beta and eta identities license replacement in all well-typed contexts. No axiom identifying currently coextensive properties is included.

Classicism $\C$ is the $\mathsf H$-theory obtained by adding all instances of
\begin{equation}\label{eq:logical-equivalence}
(\lambda\vec v.\varphi)=(\lambda\vec v.\psi)
\qquad\text{whenever }\mathsf H\vdash\varphi\leftrightarrow\psi.
\end{equation}
The abstracted variable list may be empty. This presentation fixes the logical strength used throughout the paper. It does not apply the displayed rule indiscriminately to biconditionals derived from extra non-logical axioms. The critical and set-theoretic axioms are instead explicitly adopted together with their universally closed broad necessitations.

For $\Box p=(p=\top)$, the following are the Classicist consequences used here: normal S4, converse Barcan at every type, necessary identity, and
\begin{equation}\label{eq:intensionality}
\Box\forall\vec x(F\vec x\leftrightarrow G\vec x)\rightarrow F=G.
\end{equation}
This is intensionality at every relational type; it is not current extensionality. See \citet[\S\S1.3--1.5]{BD} for its derivation and the modal properties of this calculus.

For an independently specified $\D :t\to t$, add the broadly necessitated universal closures of
$$
\D \top,\qquad \D (p\rightarrow q)\rightarrow(\D p\rightarrow \D q),\qquad
\D p\rightarrow p,\qquad \D p\rightarrow \D \D p,
$$
and converse Barcan~\textup{(9)} at every type. The basis axiom~\textup{(22)} is required for the comparison theorem, but not for the elementary results below. These $\D $ assumptions are not statements that $\D $ individuates propositions or properties.

For reference, $\HD$ denotes the weaker quantified modal fragment: $\mathsf H$ with normal S4 for $\D $ and converse Barcan, available with parameters at every $\D $-alternative. Theorems can be necessitated; undischarged assumptions cannot. Necessary identity follows from reflexivity, necessitation, and substitution using $\lambda z.\D (a=z)$. The following proofs, except the identity corollary, work in this fragment. They are written for $\D $; replacing it uniformly by $\Box$ gives their broad-rigidity analogues. In this appendix and Appendix~\ref{app:closure}, a \emph{definite} property means $\D $-rigid, as in the main text.

\subsection{The elementary results}


\begin{proof}[Proof of Proposition~3.1]
We prove the $N=\D $ case; the broad case is obtained by uniform substitution. Suppose $R$ is definite. At each alternative, inextensibility supplies the left-to-right implication in~\textup{(15)}. For the converse, assume $\D\forall x(Rx\rightarrow Qx)$. CBF gives $\forall x\D(Rx\rightarrow Qx)$. Persistence gives $Rx\rightarrow\D Rx$, so K yields $Rx\rightarrow\D Qx$. Generalize over $x$ and $Q$. The same derivation applies under the outer $\D $, since the defining clauses of definiteness supply the corresponding local premises there.

Conversely, the left-to-right implication in~\textup{(15)} is exactly inextensibility. For persistence, instantiate $Q$ with $R$ at each alternative. The theorem $\D\forall x(Rx\rightarrow Rx)$ and the reverse implication yield $\forall x(Rx\rightarrow\D Rx)$. This establishes persistence under the outer $\D $. Dualizing the biconditional, with $Q$ replaced by $\neg Q$, gives~\textup{(16)} at every alternative. This argument does not require 4.
\end{proof}

\begin{lemma}[Persistence of definiteness]\label{lem:rigrig}
$\HD$ proves $\Def(R)\rightarrow\D\Def(R)$.
\end{lemma}
\begin{proof}
Each defining conjunct is of the form $\D\varphi$. Axiom 4 and distribution over conjunction give the result.
\end{proof}

\begin{lemma}[Invariance under modal coextensiveness]\label{lem:invariance}
$R\deq S$ entails $\Def(R)\leftrightarrow\Def(S)$.
\end{lemma}
\begin{proof}
Suppose $R\deq S$ and $R$ is definite. By 4, the necessary coextensiveness is available throughout further alternatives. At such an alternative, $Sx$ implies $Rx$, and then $\D Rx$. By CBF, necessary coextensiveness gives $\D(Rx\leftrightarrow Sx)$, so $\D Sx$. This proves persistence of $S$.

For inextensibility, suppose $\forall^S x\D Qx$. Current coextensiveness gives $\forall^R x\D Qx$, whence $\D\forall^R xQx$. Necessary coextensiveness transfers this to $\D\forall^S xQx$. The argument is available at every further alternative. Symmetry proves the reverse implication.
\end{proof}

\begin{proposition}[Extensional representation]\label{prop:representation}
If $R$ and $S$ are definite, $R\sub S$ entails $\D(R\sub S)$. In particular, $R\rel S$ entails $R\deq S$.
\end{proposition}
\begin{proof}
For each current $R$ is an $S$, hence definitely an $S$ by persistence of $S$. Inextensibility of $R$, with test property $S$, gives $\D(R\sub S)$. Apply the argument in both directions for coextensiveness.
\end{proof}

\begin{corollary}\label{cor:unique}
In Classicism, every property $F$ satisfying $\ER(F)$ has exactly one broadly rigid representative, independently of the choice of $\D $.
\end{corollary}
\begin{proof}
Existence is the definition of $\ER$. Two broadly rigid representatives are coextensive, hence broadly necessarily coextensive by the broad version of Proposition~\ref{prop:representation}, hence identical by~\eqref{eq:intensionality}.
\end{proof}

The corollary is a claim of unique existence. It does not by itself supply a higher-order function which uniformly assigns representatives to arbitrary inputs. In particular, one must not infer a choice or description operator without checking its availability in the underlying intensional logic.

\begin{proposition}[Empty and singleton properties]\label{prop:singletons}
$\Z$ and $\sng a$ are definite. If a definite property is currently empty, it is definitely empty.
\end{proposition}
\begin{proof}
The empty property is persistent and every restricted universal over it is a theorem. For $\sng a$, persistence is necessary identity. If $\forall x(x=a\rightarrow\D Qx)$, reflexivity gives $\D Qa$. Necessitate the equality-substitution theorem $Qa\rightarrow\forall x(x=a\rightarrow Qx)$ and use K. The argument applies at every alternative, giving inextensibility. Finally, if $R$ is definite and empty, $\forall^R x\D\bot$ holds vacuously. Inextensibility gives $\D\forall^R x\bot$.
\end{proof}

\begin{remark}[Plural indiscernibility]
Plural variables range over $\Pl$, hence over non-empty broadly rigid properties. Corollary~\ref{cor:unique} and substitution of identicals establish their unrestricted indiscernibility. This applies even when $\D $ differs from $\Box$. The weaker result for two merely definite representatives is necessary coextensiveness under $\D $, not identity.
\end{remark}

\section{Comparing the modalities and identifying representatives}\label{app:identity}

\subsection{Transfer from broad rigidity to definiteness}

All the $\D $ assumptions of Section~3 are available at every broad alternative, and $\D $ is a term of type $t\to t$. Accordingly, $\Box p\rightarrow \D p$ follows from $\D \top$ and substitution of identicals; the same implication holds throughout broad possibility. Assume also $\mathsf{MB}$ from~\textup{(22)}.

\begin{proof}[Proof of Proposition~3.2]
Fix a broadly rigid $R$. Its broad rigidity persists under $\Box$, so it suffices to establish the local persistence and inextensibility implications at an arbitrary broad alternative. If $Rx$, broad persistence gives $\Box Rx$, hence $\D Rx$.

Now suppose $\forall^R x\,\D Qx$. Choose a basis $K$ available at this alternative. If $Kq$, the basis equation gives $\forall^R x\,\Box(q\rightarrow Qx)$. Broad inextensibility of $R$, applied to $\lambda x(q\rightarrow Qx)$, gives
$$
\Box\forall^R x(q\rightarrow Qx).
$$
The universally valid equivalence between the embedded formula and $q\rightarrow\forall^R xQx$ yields $\Box(q\rightarrow\forall^R xQx)$. Generalizing over $q$ and using the same basis gives $\D \forall^R xQx$. No selection of bases uniform across alternatives is needed.

We have established both local implications at every broad alternative. Thus
$$
\Box\forall x(Rx\rightarrow \D Rx),\qquad
\Box\forall Q(\forall^R xDQx\rightarrow \D \forall^R xQx).
$$
Replacing these outer boxes with $\D $ gives $\Def(R)$. If $\ER(F)$, take its broadly rigid representative and apply the result, obtaining $\ED(F)$.
\end{proof}

\subsection{Why normal S4 is not enough}\label{subsec:neighborhood}

The basis hypothesis cannot be replaced simply by normal S4, converse Barcan, and $\Box p\rightarrow \D p$. Here is a countermodel. Take the full action model on a frame with a root $r$ and terminal successors $s_0,s_1,\ldots$, with the reflexive arrows and the arrows $r\to s_n$. The individual action is the constant set $\mathbb N$, with identity transitions. Use the full propositional and exponential actions; this is a model of $\C$. At the root, propositions are arbitrary subsets of $\{r,s_0,s_1,\ldots\}$, and at each terminal stage they are ordinary truth values.

Define a genuine higher-order operator $\D $ by
$$
\begin{aligned}
\D  p\text{ at }r
&\quad\Longleftrightarrow\quad
p\text{ at }r\text{ and }p\text{ at all but finitely many }s_n,\\
\D  p\text{ at }s_n
&\quad\Longleftrightarrow\quad p\text{ at }s_n.
\end{aligned}
$$
The terminal values of the root operation agree with the terminal operations, so it is an element of the full exponential action, not a metalinguistic predicate. At the root it is the interior operator for neighborhoods containing the root and cofinitely many terminal points. It preserves $\top$ and finite conjunctions, is factive, and is idempotent; hence its normal S4 principles hold. Broad necessity entails $\D $-necessity.

Converse Barcan holds at every type. If $\D \forall v\varphi$ holds at the root, the universal is true there and at cofinitely many terminal stages. For each current $v$, its image satisfies $\varphi$ at all those stages, so $\D \varphi$ holds. At terminal stages the result is immediate. This reasoning applies to the expanding higher-order domains as well as the constant individual domain.

The universal individual property $\U$ is broadly rigid, because individuals do not change or increase along broad transitions. Let $Q$ be true of every $k$ at $r$ and true of $k$ at $s_n$ exactly when $k\ne n$. The full property domain contains this intension. At $r$ every $k$ is definitely $Q$, but $\forall kQk$ fails at every terminal stage and is not definite. Thus $\U$ is not $\D $-inextensible. This model refutes $\Rig(\U)\rightarrow\Def(\U)$ while satisfying the other modal assumptions. It explains why the comparison theorem explicitly includes~\textup{(22)}.

The basis condition need not collapse the modalities. In the same full action model, take instead $\D =\lambda p.p$. At each stage there is a proposition true precisely at that stage; the singleton of that proposition is a basis satisfying~\textup{(22)}. All $\D $ principles hold, but at the root many truths are not broadly necessary. This example establishes non-collapse of the comparison assumptions alone, not of their further combination with the set axioms.

\subsection{Identity of definite representatives}

The next result needs no basis assumption. Work in $\C$ expanded by a genuine normal S4 operator $\D $ satisfying the quantified assumptions. The two modalities are initially distinct.

\begin{theorem}[Representative identity forces broad necessity]\label{thm:collapse}
Principle~\textup{(27)} entails $\forall p(\D p\leftrightarrow\Box p)$.
\end{theorem}
\begin{proof}
Fix $p$ with $\D p$ and an individual $a$. Let $H=\lambda x.x=a$ and $F=\lambda x(p\wedge x=a)$. Normality gives $\D (F\rel H)$. Proposition~\ref{prop:singletons} gives $\Def(H)$, and Lemma~\ref{lem:invariance} gives $\Def(F)$. T also gives current coextensiveness. By~\textup{(27)}, $F=H$.

Substitute into application at $a$. Beta conversion yields equality, at propositional type, between $p\wedge(a=a)$ and $a=a$. In $\C$, the logical identities $a=a=\top$ and $p\wedge\top=p$ hold. Hence $p=\top$, which is $\Box p$. Conversely, if $p=\top$, substitution into the genuine operator $\D $ and the theorem $\D \top$ give $\D p$.
\end{proof}

\begin{corollary}
The principle $\D \forall x(Fx\leftrightarrow Gx)\rightarrow F=G$ also entails agreement of $\D $ with broad necessity.
\end{corollary}
\begin{proof}
Proposition~\ref{prop:representation} shows that this principle implies~\textup{(27)}.
\end{proof}

The result gives pointwise agreement at any state where the premise holds, and agreement at every state if the premise is imposed throughout the model. It does not infer identity of the operators as higher-order entities merely from agreement at one state. Broadly rigid representatives already satisfy their own identity principle; no identification of all definite representatives is needed for the plural interpretation.


\section{Closure proofs}\label{app:closure}

\begin{proposition}[Finite union]\label{prop:union}
If $R$ and $S$ are definite, so is $R\vee S:=\lambda x(Rx\vee Sx)$. Consequently $\ED$ is closed under finite union.
\end{proposition}
\begin{proof}
Persistence follows from persistence of the disjuncts and normality. At an arbitrary alternative, suppose $\forall^{R\vee S}x\D Qx$. This entails both $\forall^Rx\D Qx$ and $\forall^Sx\D Qx$. Inextensibility gives $\D\forall^RxQx$ and $\D\forall^SxQx$, hence $\D\forall^{R\vee S}xQx$. Thus the union is inextensible. For ED, instantiate its two existential witnesses and take their definite union; its extension is the union of the original extensions.
\end{proof}

\begin{proposition}[Definite separation]\label{prop:sharpsep}
If $R$ is definite and~\textup{(28)} holds, then $R\wedge F$ is definite.
\end{proposition}
\begin{proof}
Let $T(x):=(\D Fx\vee\D\neg Fx)$. By 4 and normality, $T(x)\rightarrow\D T(x)$. The hypothesis and inextensibility of $R$ therefore give $\D\forall^R xT(x)$. Hence the pointwise sharpness condition is available at every alternative; so is definiteness of $R$ by Lemma~\ref{lem:rigrig}.

At any such alternative, $Rx\wedge Fx$ entails $\D Rx$ by persistence. Sharpness and T give $Fx\rightarrow\D Fx$ on $R$, so the conjunction is persistent. For inextensibility suppose $\forall x(Rx\wedge Fx\rightarrow\D Qx)$. For any $x$ with $Rx$, if $Fx$ then $\D Qx$, hence $\D(Fx\rightarrow Qx)$. If $\neg Fx$, sharpness and T give $\D\neg Fx$, again yielding $\D(Fx\rightarrow Qx)$. Inextensibility of $R$, with $\lambda x(Fx\rightarrow Qx)$ as test, gives $\D\forall x(Rx\rightarrow(Fx\rightarrow Qx))$. This is the required restricted universal. Since the argument holds at every alternative, $R\wedge F$ is inextensible.
\end{proof}

\begin{proposition}[Negative persistence and bounded Boolean operations]\label{prop:negative}
Assume $\ND$. Every definite $S$ satisfies $\D\forall x(\neg Sx\rightarrow\D\neg Sx)$. Definite properties are closed under intersection and relative complement.
\end{proposition}
\begin{proof}
At an arbitrary alternative, fix $a$ with $\neg Sa$. Each $y$ with $Sy$ is distinct from $a$, and $\ND$ makes that distinction definite. Inextensibility of $S$, tested with $\lambda y.y\ne a$, gives $\D\forall y(Sy\rightarrow y\ne a)$. By reflexivity this implies $\D\neg Sa$. The argument applies at all alternatives. Positive and negative persistence give pointwise definiteness of $S$ and of $\neg S$. Apply Proposition~\ref{prop:sharpsep} to $R$ with these two conditions.
\end{proof}

\begin{corollary}[Generalized Boolean algebra]
Under $\ND$, definite properties modulo current coextensiveness form a generalized Boolean algebra under inclusion. If $\neg\ED(\U)$, this algebra has no greatest element. For broad necessity in Classicism the quotient is unnecessary.
\end{corollary}
\begin{proof}
The operations just established, together with $\Z$, satisfy the generalized Boolean identities extensionally, by classical logic. They are well defined on coextensiveness classes. A greatest class would contain every singleton and therefore be universal. The broad-necessity claim follows from Corollary~\ref{cor:unique}.
\end{proof}

\begin{proposition}[Uniform definite union]\label{prop:rigunion}
Equation~\textup{(31)} is a theorem of $\HD$.
\end{proposition}
\begin{proof}
Let $U(y):=\exists x(Ax\wedge Bxy)$. By Lemma~\ref{lem:rigrig}, the hypothesis about fibres yields $\forall^Ax\D\Def(B_x)$. Inextensibility of $A$ gives $\D\forall^Ax\Def(B_x)$. The indexing definiteness and the fibre definiteness are therefore available throughout the following argument.

For persistence, an instance $Uy$ has a witness $x$ with $Ax$ and $Bxy$. Persistence of $A$ and of $B_x$ gives $\D Ax$ and $\D Bxy$, whence $\D Uy$ by existential introduction inside $\D $.

For inextensibility, at an arbitrary alternative suppose $\forall^Uy\D Qy$. For each $x$ with $Ax$, we have $\forall y(Bxy\rightarrow\D Qy)$. Inextensibility of $B_x$ gives $\D\forall y(Bxy\rightarrow Qy)$. Apply inextensibility of $A$ to the test property $\lambda x\forall y(Bxy\rightarrow Qy)$. The result is $\D\forall x(Ax\rightarrow\forall y(Bxy\rightarrow Qy))$, equivalent to $\D\forall^UyQy$. This holds at every alternative.
\end{proof}

No selection of representatives occurs in this proof. In particular, it does not license replacing its premise $\Def(B_x)$ by $\ED(B_x)$.

\begin{remark}[Empty-inclusive formulations]
The empty-inclusive Generalized Union axiom used here is equivalent to its non-empty counterpart in the presence of Singleton and the resulting Separation principle. For the nontrivial direction, given a represented index extension, separate off the indices whose fibres are non-empty. If there are any such indices, apply non-empty Generalized Union to them; otherwise the union is empty and has its canonical representative. This uses no choice of representatives. Infinity's empty starting case is similarly witnessed by the empty property. Thus permitting empty representatives in the higher-order formulation does not add a further completion principle.
\end{remark}

\section{The injective level model}\label{app:model}

We work in the metatheory ZFC. All categories and type domains used below are sets. Fix a strictly increasing sequence $(\kappa_n)_{n<\omega}$ of infinite regular cardinals and identify each cardinal with its initial ordinal. Put $E_n=\kappa_n$, and set $E_{-1}=\varnothing$ for convenience. Let $\lev(a)$ be the least $j$ with $a\in E_j$.

\subsection{The category and supported action construction}

The category $\mathcal C$ has objects $W_n$, with individual domain $W_n^e=E_n$. For $n\le m$, its arrows $W_n\to W_m$ are the injections $f:E_n\to E_m$ such that
\begin{equation}\label{eq:catlevel}
f[E_j]\subseteq E_j\quad\text{for each }j\le n.
\end{equation}
For $n>m$ there are no arrows. Composition is composition of functions. The identities and the inclusions satisfy~\eqref{eq:catlevel}; composition preserves it. Thus this is a category rooted at $W_0$.

Let $\mathcal I_n=\{X\subseteq E_n:|X|<\kappa_n\}$. It contains the empty set and all singletons and is closed under subsets and finite unions. Arrows act on supports by direct image. If $f:W_n\to W_m$ and $X\in\mathcal I_n$, then $|f[X]|<\kappa_n\le\kappa_m$, so $f[X]\in\mathcal I_m$.

Here is the supported action construction used throughout this appendix. For an action $-^\sigma$ of $\mathcal C$, an element $u\in W_n^\sigma$ is \emph{pinned down by} $X\subseteq E_n$ when parallel arrows $f,g:W_n\to W_m$ agreeing on $X$ satisfy $f^\sigma u=g^\sigma u$. The full propositional action at $W$ is $\mathcal P(\Arr(W))$, with
\begin{equation}\label{eq:paction}
f^t p=\{i\in\Arr(\tg(f)):i\circ f\in p\}.
\end{equation}
Admit exactly its elements pinned down by an allowed support. At function type $\sigma\to\tau$, use the supported elements of the exponential action: families $\alpha(h,a)\in\tg(h)^\tau$, for $h$ out of $W$ and $a\in\tg(h)^\sigma$, satisfying
\begin{equation}\label{eq:naturality}
i^\tau\alpha(h,a)=\alpha(i\circ h,i^\sigma a).
\end{equation}
The action on these families is reindexing, $(f\alpha)(h,a)=\alpha(h\circ f,a)$. Type domains are obtained recursively in this way, starting from the stipulated individual action and the supported propositional action.

We use the supported-fullness theorem of \citet[Appendix~D, Proposition~D.5]{BD}: when the allowed supports contain $\varnothing$, are closed under finite unions and direct images, and pin down the individual action, this construction is an action model of Classicism. Its application here requires no cardinal-completeness assumption. Individuals have singleton supports; all the remaining hypotheses were checked above. Briefly, logical constants are supported by the empty set, and application to supported arguments is supported by the union of the transported function support and the argument support. The same mechanism validates the combinators and lambda formation.

Truth at $W$ means that $\id_W$ belongs to the proposition denoted. At a stage reached by an arrow $h$, quantifiers range over the domains of $\tg(h)$ and connectives are classical. Equality is equality in these domains. Broad necessity $p=\top$ holds just when $p$ contains every outgoing arrow. In particular,
\begin{equation}\label{eq:modaltruth}
\mathcal A,h,g\models\Box\varphi
\quad\Longleftrightarrow\quad
\mathcal A,i\circ h,i\circ g\models\varphi
\ \text{for every }i\in\Arr(\tg(h)).
\end{equation}
In this model we set $\D =\Box$. Every individual transition is injective, so $\ND$ holds everywhere.

\subsection{Unary intensions and hulls}

Let
$$
P_W=\{(h,a):h\in\Arr(W),\ a\in\tg(h)^e\}.
$$
For a property $\alpha\in W^{e\to t}$ define its intension by
$$
I_\alpha=\{(h,a):\id_{\tg(h)}\in\alpha(h,a)\}.
$$
Write $I_\alpha(h)=\{a:(h,a)\in I_\alpha\}$ for its extension at $h$.

\begin{lemma}[Intension representation]\label{lem:intensions}
The unary properties at $W_n$ correspond bijectively to the sets $I\subseteq P_{W_n}$ which are $X$-invariant for some $X\in\mathcal I_n$, where $X$-invariance means that
$$
(h,a)\in I\ \leftrightarrow\ (h',a)\in I
$$
whenever $h,h'$ are parallel and agree on $X$. The inverse correspondence is
\begin{equation}\label{eq:inverse}
\alpha_I(h,a)=\{i\in\Arr(\tg(h)):(i\circ h,i(a))\in I\}.
\end{equation}
Moreover,
\begin{equation}\label{eq:transportI}
I_{f\alpha}(h)=I_\alpha(h\circ f).
\end{equation}
\end{lemma}
\begin{proof}
Equation~\eqref{eq:naturality} implies that a property's values are recovered from its intension by~\eqref{eq:inverse}. Conversely, that equation defines a natural family. If $I$ is $X$-invariant, its value at $(h,a)$ is a proposition pinned down by $h[X]\cup\{a\}$: arrows agreeing on this support have the same action on $a$ and agree after $h$ on $X$, including after further composition. The family is pinned down by $X$. Conversely, a family pinned down by $X$ gives $X$-invariance by evaluating the equality of transported families at an identity arrow and at $a$. Thus both constructions land in the admitted domains and are inverse. The transport equation is immediate from reindexing.
\end{proof}

The same argument applies to binary relations, using $h[X]\cup\{a,b\}$ for the support of a value. In particular, any subset of the current individual domain is the extension of a property: for $Y\subseteq E_n$, let $(h,a)\in I$ just when $a\in Y$. This is empty-supported as an intension, regardless of the size of $Y$. The analogous observation applies to arbitrary current binary relations. This is \emph{extensional} fullness at these types, not the availability of every rigid intension.

For $Y\subseteq W^e$, define the hull intension
\begin{equation}\label{eq:hull}
\Hull_W(Y)=\{(h,a):a\in h[Y]\}.
\end{equation}
When $Y$ is small, it is $Y$-invariant, so it represents a property $\alpha_Y$.

\begin{lemma}[Hull properties]\label{lem:hulls}
For small $Y$ at $W$, $\alpha_Y$ is rigid, has current extension $Y$, and satisfies $f\alpha_Y=\alpha_{f[Y]}$. If a rigid property has small current extension $Y$, it is $\alpha_Y$.
\end{lemma}
\begin{proof}
The current extension and transport equation follow directly from~\eqref{eq:hull}. Persistence follows because $a\in h[Y]$ implies $i(a)\in(i\circ h)[Y]$. For inextensibility at the stage after $h$, suppose each $a\in h[Y]$ is definitely $Q$. At every further arrow $i$, all the objects of $i[h[Y]]$ satisfy the transported $Q$. These are exactly the instances of the transported hull, giving the required definite restricted universal. The argument applies at every stage. Finally, a rigid property with extension $Y$ is coextensive with the rigid $\alpha_Y$ and is identical to it by Corollary~\ref{cor:unique}.
\end{proof}

\subsection{Classification}

For each $n$, let $\rho_n\in W_n^{e\to t}$ have intension
$$
I_{\rho_n}=\{(h,a):a\in E_n\}.
$$
It is empty-supported. Every current individual $a\in E_n$ definitely satisfies $\rho_n$, since every arrow out of $W_n$ maps $E_n$ into itself.

\begin{theorem}[Classification of rigid properties]\label{thm:classification}
At every $W_n$, the rigid properties are exactly the hulls $\alpha_Y$ for $Y\subseteq E_n$ of size $<\kappa_n$.
\end{theorem}
\begin{proof}
Small hulls are rigid by Lemma~\ref{lem:hulls}. Conversely, let $\alpha$ be rigid, pinned down by $X\in\mathcal I_n$, with current extension $Y$. If $Y$ is small, the same lemma identifies $\alpha$ with $\alpha_Y$.

Suppose instead that $|Y|=\kappa_n$. Since every current individual definitely satisfies $\rho_n$, inextensibility gives
$$
\Box\forall x(\alpha x\rightarrow\rho_nx).
$$
Let $j:W_n\to W_{n+1}$ be the inclusion and let $S$ be the current extension of $j\alpha$. The displayed formula gives $S\subseteq E_n$. Persistence along $j$ gives $Y\subseteq S$. Since $|E_n|=\kappa_n<\kappa_{n+1}$, $S$ is small at $W_{n+1}$. The transported property remains rigid and thus, by Lemma~\ref{lem:hulls},
$$
j\alpha=\alpha_S.
$$
It is pinned down by $X$: parallel arrows agreeing on $X=j[X]$ yield composites with $j$ agreeing on the original support.

Choose $b\in Y\setminus(E_{n-1}\cup X)$. Such a choice is possible because $|Y|=\kappa_n$, whereas $|E_{n-1}\cup X|<\kappa_n$. Choose distinct $b=c_0,c_1,c_2,\ldots$ in $E_n\setminus(E_{n-1}\cup X)$. Define $k:E_{n+1}\to E_{n+1}$ by $k(c_r)=c_{r+1}$ and by fixing all other elements. This is an injective, level-respecting self-arrow of $W_{n+1}$. It agrees with the identity on $X$ and omits $b$ from its range.

Support invariance gives $k(j\alpha)=j\alpha$. But hull transport gives $k(j\alpha)=k\alpha_S=\alpha_{k[S]}$. Their current extensions would therefore satisfy $S=k[S]$, impossible since $b\in Y\subseteq S$ and $b\notin k[S]$. This excludes a large extension and completes the classification.
\end{proof}

The classification argument uses the increasing cardinal bounds and the availability of an infinite sequence outside a small support. Regularity is needed for the closure claims below, not for this part of the proof.

\subsection{Critical plural logic and Infinity}

Translate $x\prec xx$ as $Xx$, and translate plural quantifiers by restricting property quantifiers to $\Pl(X)$. Coextensiveness of pluralities is translated by $X\rel Y$. Non-emptiness is built in. Indiscernibility follows from identity of coextensive rigid properties in $\C$; thus it applies even to formulas with additional higher-order parameters and modal operators. Singleton and Adjunction follow from the logical results already proved.

\begin{theorem}[Validation of the critical principles]\label{thm:critical}
For any increasing sequence of infinite regular cardinals, Generalized Union, Separation, and Replacement hold at every stage of the model, under every assignment. If every $\kappa_n$ is uncountable, Infinity also holds at every stage. There is no universal plurality at any stage, and $\ND$ holds at every stage.
\end{theorem}
\begin{proof}
Fix a stage with individual domain $E_n$ and an assignment. By Theorem~\ref{thm:classification}, rigid properties have exactly the small subsets as their extensions, and non-empty rigid properties have exactly the non-empty small subsets. Interpret a formula $\psi(x,y)$ under this assignment by its current relation $\Psi\subseteq E_n^2$, writing $\Psi_a=\{b:(a,b)\in\Psi\}$. This is legitimate for arbitrary formulas, including those with modal and higher-order quantifiers: their satisfaction relation still determines a subset of $E_n^2$.

For Generalized Union let $A\subseteq E_n$ be small and suppose every $\Psi_a$, for $a\in A$, is small. Regularity of $\kappa_n$ gives $|\bigcup_{a\in A}\Psi_a|<\kappa_n$. Its hull is the required representative. If the plural premises require $A$ and its fibres to be non-empty, their union is non-empty as well. Separation follows because a subset of a small set is small. Replacement follows because the image of a small set under a function has at most the same cardinality.

For Infinity suppose $\Psi$ is the graph of a total function $f:E_n\to E_n$ and let $Y$ be small. Put $C=\bigcup_{r<\omega}f^r[Y]$, including $f^0[Y]=Y$. Then $C$ contains $Y$, is closed under $f$, and has cardinality at most $\max(\aleph_0,|Y|)$ when $Y$ is non-empty. This is $<\kappa_n$ if $\kappa_n$ is uncountable. The empty case is immediate. The hull of $C$ witnesses the principle, and is non-empty when required.

There is no universal plurality because $E_n$ is not small. Injectivity of all transitions establishes necessary distinctness of individuals. Since the argument applies to every stage and every assignment, the principles hold after every arrow from the root and so do their necessitations.
\end{proof}

\begin{corollary}[A ZFC construction]\label{cor:modelzfc}
Taking $\kappa_n=\aleph_{n+1}$ gives a model of $\C$ in which all principles of critical plural logic, $\ND$, and the absence of a universal plurality hold necessarily. In particular, rigid comprehension fails at individual-property type. No inaccessible-cardinal assumption is used.
\end{corollary}

\begin{corollary}[Independence of Infinity]\label{cor:noinfinity}
Taking $\kappa_n=\aleph_n$ validates the other principles just listed, but refutes Infinity at the root.
\end{corollary}
\begin{proof}
The root individual domain is $\omega$, and its rigid extensions are exactly finite subsets. By Lemma~\ref{lem:intensions} and its binary analogue, there is a relation whose current extension is $\{(n,n+1):n\in\omega\}$. At the root it is total and functional. The singleton of $0$ has no finite superset closed under that relation. Thus the plural Infinity instance for this relation fails. The remaining claims follow from Theorem~\ref{thm:critical}.
\end{proof}

\section{The set-theoretic derivations and their joint model}\label{app:sets}

\subsection{Deriving the set axioms}

We work in $\mathsf{CS}$ as defined in Section~7. Every axiom below is available with higher-order and modal parameters. A set is an individual satisfying $\mathsf S$, and $E_a=\lambda x.x\in a$. Equations~\textup{(38)},~\textup{(39)} and Proposition~3.2 give
$$
\ER(F)\leftrightarrow\ED(F)\leftrightarrow
\exists a(\mathsf Sa\wedge E_a\rel F).
$$
Uniqueness of the associated set follows from Extensionality. It will be used only as uniqueness of a witness to an existential claim; no uniform choice operator is inferred.

\begin{proof}[Proof of Theorem~7.1]
\emph{Empty Set and Pairing.} The broadly rigid properties $\Z$ and $\lambda x(x=a\vee x=b)$ have associated sets. Extensionality makes the empty set unique. Denote it informally by $0$ in the remainder of the proof, with existential elimination understood.

\emph{Separation.} Given a set $a$ and any formula $\varphi(x)$, let $F=\lambda x(x\in a\wedge\varphi(x))$. Set-boundedness gives $\ER(E_a)$; broad Separation gives $\ER(F)$. Its associated set $b$ satisfies $x\in b\leftrightarrow(x\in a\wedge\varphi(x))$. This proves the full expanded-language schema, not just a first-order sublanguage instance.

\emph{Replacement.} Suppose $\forall x(x\in a\rightarrow\exists!y\psi(x,y))$. Each fibre over a member of $a$ has a singleton rigid representative. Generalized Union gives a rigid representative of $\lambda y\exists x(x\in a\wedge\psi(x,y))$. Its associated set is the required image. If $a$ is empty, the empty representative gives the same conclusion.

\emph{Union.} For each $x\in a$, $E_x$ has a rigid representative: by set-boundedness if $\mathsf Sx$, and by the empty property if $\neg\mathsf Sx$, since only sets have members. Generalized Union gives a rigid representative of $\lambda y\exists x(x\in a\wedge y\in x)$. Its associated set is the union of $a$.

\emph{Infinity.} Define a formula $\theta(x,y)$ by
$$
\begin{aligned}
\theta(x,y)\quad\leftrightarrow\quad
&\bigl(\mathsf Sx\wedge\mathsf Sy\wedge
       \forall z(z\in y\leftrightarrow(z\in x\vee z=x))\bigr)\\
&\vee\bigl(\neg\mathsf Sx\wedge y=0\bigr).
\end{aligned}
$$
For sets, Pairing and Union provide the indicated successor; alternatively use the union of $E_x$ and its singleton representative and then Set Formation. Extensionality gives uniqueness. For urelements the default value is unique. Thus $\forall x\exists!y\theta(x,y)$.

Apply broad Infinity to the rigid singleton $\sng 0$ and $\theta$. Obtain a broadly rigid $R$ containing $0$ and closed under $\theta$. Associate a set $c$ with $R$ and use Separation to form the subset $i$ of its set members. Then $0\in i$. If $x\in i$, its unique successor $y=x\cup\{x\}$ belongs to $R$ by closure and is a set, so $y\in i$. This is the usual inductive-set axiom, with the members of $i$ explicitly restricted to sets.

Extensionality is an axiom. These constructions use formulas and existential witnesses only, so they introduce neither a property-to-set function nor a choice principle.
\end{proof}

\begin{proposition}[No universal represented extension]
The theory $\mathsf{CS}$ proves $\neg\ER(\U)$ and $\neg\ED(\U)$. It also proves $\neg\ER(\mathsf S)$ and $\neg\ED(\mathsf S)$.
\end{proposition}
\begin{proof}
If $\ER(\U)$, Separation gives $\ER(\lambda x.x\notin x)$; its associated set yields Russell's contradiction. If $\ER(\mathsf S)$, use instead $\lambda x(\mathsf Sx\wedge x\notin x)$. Its associated individual $r$ is itself a set, so again $r\in r\leftrightarrow r\notin r$. The claims about $\ED$ follow from the displayed equivalence.
\end{proof}

\subsection{The cardinal preparation}

For a regular uncountable cardinal $\kappa$, let
$$
H_\kappa=\{x:|\operatorname{TC}(\{x\})|<\kappa\}.
$$
The transitive-closure convention includes the object itself and does not affect the usual definition. The following elementary facts will be used explicitly.

\begin{lemma}\label{lem:hk}
Suppose $\kappa$ is regular and uncountable and $2^{<\kappa}=\kappa$. Then:
\begin{enumerate}[label=(\roman*)]
\item $H_\kappa$ is transitive and has cardinality $\kappa$.
\item Every member of $H_\kappa$ has cardinality below $\kappa$.
\item If $Y\subseteq H_\kappa$ and $|Y|<\kappa$, then $Y\in H_\kappa$.
\end{enumerate}
\end{lemma}
\begin{proof}
Transitivity and (ii) follow from the definition. For (iii), the transitive closure of $\{Y\}$ is contained in
$$
\{Y\}\cup\bigcup_{y\in Y}\operatorname{TC}(\{y\}).
$$
Regularity makes this union small. To bound $|H_\kappa|$, code each element by a pointed well-founded extensional relation on some cardinal $\mu<\kappa$, representing its transitive closure. There are at most $2^\mu\cdot\mu$ such codes for infinite $\mu$, and only finitely or countably many of the finite sizes. The sum over $\mu<\kappa$ is at most $\kappa$ under $2^{<\kappa}=\kappa$. Conversely, all ordinals below $\kappa$ belong to $H_\kappa$, so its cardinality is exactly $\kappa$.
\end{proof}

Work now in ZFC+GCH and take $\kappa_n=\aleph_{n+1}$. These cardinals are regular and satisfy the lemma's cardinal equation. Choose bijections
$$
b_n:E_n\longrightarrow H_{\kappa_n},\qquad E_n=\kappa_n.
$$
There is no requirement that these bijections commute with the individual transitions. That would impose a different, modal condition on membership.

\subsection{Membership as an admitted intension}

Use exactly the category, individual action, support ideals, and higher-order domains of Appendix~\ref{app:model}. For a binary relation at $W_n$, the analogue of Lemma~\ref{lem:intensions} identifies admitted intensions with suitably supported subsets of triples $(h,a,b)$, where $h:W_n\to W_m$ and $a,b\in E_m$. Define
\begin{equation}\label{eq:membership-intension}
I^n_{\in}(h,a,b)
\quad\Longleftrightarrow\quad b_m(a)\in b_m(b).
\end{equation}
The membership symbol on the right is the metatheoretic membership relation. The left side specifies an object-language binary intension. It depends on the target stage and the arguments, but not on the particular arrow to that stage. Thus it is invariant under all parallel arrows and has empty support. The corresponding value at $(h,a,b)$ is a proposition supported by $\{a,b\}$, so the binary analogue of the representation lemma proves that this relation belongs to the higher-order domain.

Moreover, transport is coherent. If $f:W_n\to W_m$, the transported intension at $(h,a,b)$ is $I^n_{\in}(h\circ f,a,b)$, which equals $I^m_{\in}(h,a,b)$. Hence these stagewise intensions interpret one persistent language constant, with its ordinary action-model transport. This is the needed distinction between constructing one higher-order model and merely assigning a different first-order structure at each node. Coherence of the \emph{relation's transport} does not imply preservation of the \emph{truth} of each membership application along individual transitions.

Interpret $\mathsf S$ by $\U$ and $\D $ by $\Box$. The latter satisfies all modal assumptions and the basis condition, witnessed by $Kq:=(q=\top)$. At $W_n$ the current membership structure is isomorphic, via $b_n$, to $(H_{\kappa_n},\in)$.

\subsection{Verification and relative consistency}

\begin{proof}[Proof of Theorem~7.2]
The underlying action model validates $\C$, necessary distinctness at individual type, and all the broad critical principles by Theorem~\ref{thm:critical}. Since $\D =\Box$, $\ED(F)$ and $\ER(F)$ both hold exactly when the current extension of $F$ has size below $\kappa_n$.

Typing is immediate, as every individual is a set. Extensionality follows from extensionality in the transitive structure $H_{\kappa_n}$. If $s$ is a current individual, its extension under membership corresponds to the elements of $b_n(s)$, a small set by Lemma~\ref{lem:hk}(ii). It therefore has a hull representative. This proves set-boundedness.

For Set Formation, let $F$ have small current extension $Y\subseteq E_n$. Its image $b_n[Y]$ is a small subset of $H_{\kappa_n}$ and so, by Lemma~\ref{lem:hk}(iii), is itself a member of $H_{\kappa_n}$. Choose the label $s=b_n^{-1}(b_n[Y])$. Then, for every current $x$,
$$
x\in s\quad\Longleftrightarrow\quad
b_n(x)\in b_n[Y]\quad\Longleftrightarrow\quad x\in Y.
$$
Thus $s$ represents the current extension of $F$. This proof applies to every admitted $F$, not just properties definable without parameters. In particular, it handles conditions containing membership, higher-order quantifiers, and modal operators. It also applies at every stage and every assignment. Coherence of the membership intension therefore makes the broad necessitations of all the axioms true.

Foundation holds in each transitive $H_{\kappa_n}$. Even the higher-order minimal-element formulation holds, since each non-empty current property extension is a metatheoretic set and has an $\in$-minimal member. Choice holds as well: given a current set of non-empty sets, metatheoretic Choice selects one member of each; the resulting function, coded by ordered pairs, has transitive closure of size below $\kappa_n$, by regularity. It belongs to $H_{\kappa_n}$. Both principles therefore hold at every broad alternative.

To refute Powerset at $W_n$, let $a=\aleph_n$, regarded as an element of $H_{\aleph_{n+1}}$. Every subset of $a$ belongs to that structure. But the collection of all such subsets has size $2^{\aleph_n}=\aleph_{n+1}$ and is not a member of $H_{\aleph_{n+1}}$. No smaller internal candidate can have precisely all those subsets as members. Powerset fails at each stage.

The construction was carried out in ZFC+GCH. If ZFC has a model, its constructible universe satisfies ZFC+GCH and carries the construction. Equivalently, the argument is formalizable there as a relative consistency proof. No inaccessible cardinal, or set-sized model of ZFC itself as an individual domain, is assumed.
\end{proof}

The extra GCH assumption used during construction should not be erased from the proof: the expansion needs enough individual labels to represent every small subset. Relative consistency with ZFC follows by passing to a model of GCH, rather than by asserting that $|H_{\aleph_{n+1}}|=\aleph_{n+1}$ holds in every universe of ZFC. The argument establishes the existence of the combined model, and not an expansion of every choice of the original cardinal data.


\section{Countermodels and limitations}\label{app:counter}

\subsection{Separation, Replacement, and Generalized Union need more}

Apply the supported action construction to the one-object category whose arrows are all permutations of $\omega$, acting on individuals in the evident way. Allow finite supports. This is a model of $\C$; all transitions are invertible and the broad modality is S5.

\begin{lemma}
The rigid extensions in this model are exactly the finite and cofinite subsets of $\omega$.
\end{lemma}
\begin{proof}
If a property $\alpha$ is persistent and has current extension $Y$, then $I_\alpha(h)=h[Y]$ for every permutation $h$. One inclusion follows by persistence from the identity to $h$. For the other, transport an instance at $h$ back along $h^{-1}$ and use persistence again. Thus every rigid property has the hull intension of its extension.

A finite $Y$ supports its own hull. A cofinite $Y$ has its hull supported by the finite complement, since permutations are bijective. These hulls are rigid by the same direct argument as Lemma~\ref{lem:hulls}, which requires only that the hull belong to the domain. If both $Y$ and its complement are infinite, then for any proposed finite support $X$ there are $y\in Y\setminus X$ and $z\notin Y\cup X$. The transposition exchanging $y,z$ fixes $X$ but changes $Y$, so the hull is not pinned down by $X$. It is not admitted to the domain.
\end{proof}

Let $E$ be the even natural numbers. The qualitative intension $I_F(h,a)\leftrightarrow a\in E$ is empty-supported and hence is an admitted property. Its extension $E$ has no rigid representative. Since the universal property is rigid, unrestricted Separation applied to it and $F$ fails.

Likewise, let $B$ have the qualitative extension $Bxy\leftrightarrow y=2x$ at each arrow. This relation is admitted by the binary intension construction. Every current fibre has a singleton rigid representative. With the rigid universe as indexing property, the union has extension $E$ and is not ED. Thus Generalized Union and Replacement fail as well. There is no contradiction with Proposition~\ref{prop:rigunion}: the actual fibre properties supplied by $B$ are not a uniform rigid family. The countermodel establishes non-derivability of the stronger closure principles from the base analysis, even with definite individual distinctness. It does not claim their independence after every other critical axiom is added.

\subsection{The constant-domain and S5 obstructions}

For the $\D $ interpretation in standard constant-domain Kripke semantics, $\forall x\D Qx\rightarrow\D\forall xQx$ holds because both sides quantify over the same individual domain and the same accessible alternatives. Hence $\U$ is definite. If each subset of that domain additionally has a constant-extension property in the relevant property domain, then every current property extension has a definite representative. The second conclusion uses fullness and does not follow from constant domains alone.

There is also a proof-theoretic obstruction. In the present classical quantified framework, assume S5 in addition to CBF. If $\forall x\D Qx$, then by axiom B,
$$
\D\poss\forall x\D Qx.
$$
The dual of CBF gives $\poss\forall x\D Qx\rightarrow\forall x\poss\D Qx$. Necessitating this and using K gives $\D\forall x\poss\D Qx$. S5 proves $\poss\D Qx\rightarrow Qx$, so $\D\forall xQx$. Thus Barcan holds; applying the argument at every alternative gives the necessitated form. By~\textup{(24)}, $\U$ is definite. If $\D =\Box$, there is consequently a universal plurality. For an independent $\D $, that last inference requires an additional bridge, such as~\textup{(40)}; the modal argument alone establishes $\ED(\U)$.

\subsection{The different definiteness claims really separate}

In the injective model of Appendix~\ref{app:model}, $\ID_\D (\U)$ holds and $\ED(\U)$ fails. For the opposite separation, there are propositions $p$ with $p\wedge\neg\D p$: fix distinct $a,b$ in the initial domain and use the proposition whose intension at $h$ says $h(a)=a$. It is supported by $\{a\}$, true at the identity, and false after the transposition of $a,b$. The property $F_p$ of~\textup{(25)} is ED but not ID, as shown in Section~4.

ED itself need not persist. With such a $p$, put
$$
G:=\lambda x\bigl((p\wedge x=a)\vee\neg p\bigr).
$$
At the root $G$ is coextensive with a singleton and is ED. After an arrow making $p$ false it is universal and is not ED. This is consistent with the persistence of rigidity: the original property $G$ is a changing selector, not its own rigid representative.

\subsection{What has not been claimed}

The classification theorem concerns broad rigidity at $e\to t$, not all higher-type rigid properties. The set expansion supplies an admitted membership relation and a current correspondence between sets and represented extensions. It does not supply a rigid set-forming operation or establish essential membership; indeed, membership truth can change along its transitions. Powerset is false in that expansion. Neither a completeness theorem for the critical theory nor a reduction of ordinary vagueness has been claimed. Finally, the model of the combined set theory takes $\D =\Box$; the separate comparison assumptions permit distinct modalities, but their joint consistency with all the set axioms and $\D \ne\Box$ has not been established here.


\section*{Acknowledgment and provenance}
The motivating proposal and a preliminary level-model construction were supplied by Andrew Bacon. This draft was generated with ChatGPT and revised in response to his annotations of 30--31 October 2026 and shortened for the competition on 31 October 2026. The injective construction and its classification argument develop that preliminary model; the set expansion and the comparison condition for the two modalities are additional developments in the revision.

\begin{thebibliography}{99}
\bibitem[Bacon(2025)]{BaconMath}
Bacon, Andrew. 2025. \emph{Could the truths of mathematics have been different?} Manuscript dated 9 October 2025. \url{https://andrew-bacon.github.io/papers/Could%20mathematics%20have%20been%20different.pdf}.

\bibitem[Bacon and Dorr(2024)]{BD}
Bacon, Andrew, and Cian Dorr. 2024. ``Classicism.'' In \emph{Higher-Order Metaphysics}, edited by Peter Fritz and Nicholas K. Jones, 109--190. Oxford: Oxford University Press. Section and appendix references follow the authors' online draft. \url{https://andrew-bacon.github.io/papers/Classicism.pdf}.

\bibitem[Crosilla and Linnebo(2024a)]{CLDef}
Crosilla, Laura, and \O ystein Linnebo. 2024a. ``Definiteness in Early Set Theory.'' \emph{Journal for the Philosophy of Mathematics} 1: 43--62. \url{https://doi.org/10.36253/jpm-2933}.

\bibitem[Crosilla and Linnebo(2024b)]{CLWeyl}
Crosilla, Laura, and \O ystein Linnebo. 2024b. ``Weyl and Two Kinds of Potential Domains.'' \emph{No\^us} 58: 409--430. First published online in 2023. \url{https://doi.org/10.1111/nous.12457}.

\bibitem[Fine(1975)]{Fine}
Fine, Kit. 1975. ``Vagueness, Truth and Logic.'' \emph{Synthese} 30: 265--300. \url{https://doi.org/10.1007/BF00485047}.

\bibitem[Florio and Linnebo(2021)]{FL}
Florio, Salvatore, and \O ystein Linnebo. 2021. \emph{The Many and the One: A Philosophical Study of Plural Logic}. Oxford: Oxford University Press.

\bibitem[Linnebo(2013)]{Linnebo2013}
Linnebo, \O ystein. 2013. ``The Potential Hierarchy of Sets.'' \emph{The Review of Symbolic Logic} 6: 205--228. \url{https://doi.org/10.1017/S1755020313000014}.

\bibitem[Linnebo(2026)]{Linnebo2026}
Linnebo, \O ystein. 2026. \emph{Strict potentialism in modal mirrors}. arXiv:2604.17296v1, 19 April 2026; manuscript version dated 6 March 2026. \url{https://arxiv.org/abs/2604.17296}.

\bibitem[Roberts(2022)]{Roberts}
Roberts, Sam. 2022. ``Pluralities as Nothing Over and Above.'' \emph{The Journal of Philosophy} 119(8): 405--424. \url{https://doi.org/10.5840/jphil2022119828}.

\bibitem[Williamson(1999)]{Williamson}
Williamson, Timothy. 1999. ``On the Structure of Higher-Order Vagueness.'' \emph{Mind} 108(429): 127--143. \url{https://doi.org/10.1093/mind/108.429.127}.
\end{thebibliography}
\end{document}
